\begin{lemma}[The sup-convolution duality for a separated quadratic test]
Let $E$ be a finite-dimensional real inner-product space, let $N$ be a natural
number and write $V=\ell^2\bigl(\bigoplus_{j=1}^NE\bigr)$ for the
$\ell^2$-product of $N$ copies of $E$, so that $|W|^2=\sum_j|w_j|^2$ and
$\langle Z,W\rangle=\sum_j\langle z_j,w_j\rangle$. Let $C\ge0$ and let
$G_j:E\to\R$ $(1\le j\le N)$ be continuous with $G_j(w)\le C$ for every $w$ and
$G_j(0)=0$. Let $B,A_\varepsilon:V\to V$ be self-adjoint continuous linear
operators, let $\varepsilon>0$, and assume
\begin{align}
  A_\varepsilon\xi&=B\bigl(\xi+\varepsilon A_\varepsilon\xi\bigr)
    &&(\xi\in V),\label{eq:srdt-res}\\
  \bigl(\xi+\varepsilon A_\varepsilon\xi\bigr)
    -\varepsilon B\bigl(\xi+\varepsilon A_\varepsilon\xi\bigr)&=\xi
    &&(\xi\in V),\label{eq:srdt-inv}\\
  \varepsilon^{-1}\langle\xi,\xi\rangle+\langle A_\varepsilon\xi,\xi\rangle&>0
    &&(\xi\in V,\ \xi\ne0).\label{eq:srdt-pos}
\end{align}
Assume finally the strict quadratic test inequality
\begin{equation}\label{eq:srdt-G3}
  \sum_jG_j(w_j)-\tfrac12\langle B W,W\rangle\le-\sum_j|w_j|^4
  \qquad(W\in V,\ W\ne0).
\end{equation}
Write $\mathcal G(W)=\sum_jG_j(w_j)$ and let $\mathcal G^\varepsilon$ and
$G_j^\varepsilon$ denote the sup-convolutions of \noderef{SupConvolution} with
parameter $\varepsilon$, formed in $V$ and in $E$ respectively. Then
\begin{equation}\label{eq:srdt-sep}
  \mathcal G^\varepsilon(Z)=\sum_jG_j^\varepsilon(z_j)\qquad(Z\in V),
\end{equation}
and the function
$\Upsilon(Z)=\mathcal G^\varepsilon(Z)
-\tfrac12\langle A_\varepsilon Z,Z\rangle$ satisfies
\[
  \Upsilon(0)=0,\qquad\Upsilon(Z)<0\ \ (Z\ne0),
\]
is $\lambda$-semiconvex on $V$ in the sense of \noderef{SemiconvexOn} with
$\lambda=\varepsilon^{-1}+\|A_\varepsilon\|$, and is continuous on $V$.
\end{lemma}
\begin{proof}
Throughout, $\mathcal G$ is continuous on $V$ (a finite sum of continuous
functions of the independent coordinates) and $\mathcal G\le NC$, hence
$\mathcal G(W)\le NC+|W|$; likewise $G_j(w)\le C+|w|$. By part (1) of
\noderef{SupConvolutionSemiconvex}, applied in $V$ to $\mathcal G$ with growth
constant $NC$ and in $E$ to each $G_j$ with growth constant $C$, every defining
set
$\{\mathcal G(W)-|Z-W|^2/(2\varepsilon):W\in V\}$, resp.
$\{G_j(w)-|z-w|^2/(2\varepsilon):w\in E\}$, is bounded above; so
$\mathcal G^\varepsilon$ and the $G_j^\varepsilon$ are finite (and are genuine
suprema, not the junk value of an unbounded set).

\emph{Step 1: separability \eqref{eq:srdt-sep}.}
First dispose of the degenerate case $N=0$. Then $V$ is the zero space, the
index set is empty, $\mathcal G(W)=0$ for the single element $W=0$ of $V$, and
by the defining formula of \noderef{SupConvolution}
$\mathcal G^\varepsilon(Z)=\sup_{W\in V}\bigl(\mathcal G(W)
-|Z-W|^2/(2\varepsilon)\bigr)=0-0=0$ for the single element $Z=0$ of $V$,
while $\sum_jG_j^\varepsilon(z_j)$ is an empty sum, also $0$; so
\eqref{eq:srdt-sep} holds. (The remaining clauses of the lemma are then
immediate: $A_\varepsilon$ is the zero operator on the zero space, so
$\Upsilon\equiv0$ on $V=\{0\}$, which gives $\Upsilon(0)=0$, makes the clause
``$\Upsilon(Z)<0$ for $Z\ne0$'' vacuous, makes $\Upsilon$ continuous, and makes
$\Upsilon+\tfrac\lambda2|\cdot|^2\equiv0$ convex on $V$, i.e.\ $\Upsilon$
$\lambda$-semiconvex in the sense of \noderef{SemiconvexOn}.) So assume
$N\ge1$ for the rest of Step 1 and for Steps 2--4.

Since $|Z-W|^2=\sum_j|z_j-w_j|^2$, the defining formula of
\noderef{SupConvolution} gives
\[
  \mathcal G^\varepsilon(Z)
   =\sup_{W\in V}\sum_j\Bigl(G_j(w_j)-\frac{|z_j-w_j|^2}{2\varepsilon}\Bigr)
   =\sum_j\sup_{w_j\in E}\Bigl(G_j(w_j)-\frac{|z_j-w_j|^2}{2\varepsilon}\Bigr)
   =\sum_jG_j^\varepsilon(z_j).
\]
The middle equality holds because the expression is a sum of $N\ge1$ functions
of the independent variables $w_1,\dots,w_N$, each bounded above by the previous
paragraph: the supremum of the sum is at most the sum of the suprema since each
summand is at most its own supremum, and it is at least the sum of the suprema
because for every $\theta>0$ one may choose, for each $j$, a $w_j$ with
$G_j(w_j)-|z_j-w_j|^2/(2\varepsilon)>G_j^\varepsilon(z_j)-\theta/N$ (legitimate
since $N\ge1$, so $\theta/N>0$, and $G_j^\varepsilon(z_j)$ is a genuine
supremum) and add the $N$ inequalities, obtaining a value
$>\sum_jG_j^\varepsilon(z_j)-\theta$; as $\theta>0$ was arbitrary, the
supremum of the sum is at least $\sum_jG_j^\varepsilon(z_j)$.

\emph{Step 2: the duality identity.}
For $(Z,W)\in V\times V$ put
\[
  F(Z,W)=\mathcal G(W)-\frac{|Z-W|^2}{2\varepsilon}
    -\tfrac12\langle A_\varepsilon Z,Z\rangle ,
\]
so that $\Upsilon(Z)=\sup_WF(Z,W)$ by the defining formula of
\noderef{SupConvolution} applied to $\mathcal G$ in $V$.

\emph{(2a) The two operators and the two-sided identity.}
Put
\[
  S\xi=\xi+\varepsilon A_\varepsilon\xi,
  \qquad
  T\xi=\xi-\varepsilon B\xi\qquad(\xi\in V),
\]
two continuous linear operators on $V$. Both are self-adjoint: for all
$x,y\in V$,
\[
  \langle Sx,y\rangle=\langle x,y\rangle
    +\varepsilon\langle A_\varepsilon x,y\rangle
   =\langle x,y\rangle+\varepsilon\langle x,A_\varepsilon y\rangle
   =\langle x,Sy\rangle,
\]
by the self-adjointness of $A_\varepsilon$ assumed in the statement, and
likewise $\langle Tx,y\rangle=\langle x,Ty\rangle$ by the self-adjointness of
$B$. Hypothesis \eqref{eq:srdt-inv} says exactly
\begin{equation}\label{eq:srdt-TS}
  T(S\xi)=\xi\qquad(\xi\in V),
\end{equation}
i.e.\ $TS=\mathrm{id}$. We claim the identity also holds in the other order.
Indeed, for all $x,y\in V$,
\[
  \langle S(Tx),y\rangle
   =\langle Tx,Sy\rangle
   =\langle x,T(Sy)\rangle
   =\langle x,y\rangle,
\]
the first equality by the self-adjointness of $S$, the second by the
self-adjointness of $T$, and the third by \eqref{eq:srdt-TS} applied to
$\xi=y$. Since $\langle S(Tx)-x,y\rangle=0$ for every $y\in V$, taking
$y=S(Tx)-x$ gives $|S(Tx)-x|^2=0$, so
\begin{equation}\label{eq:srdt-ST}
  S(Tx)=x\qquad(x\in V),
\end{equation}
i.e.\ $ST=\mathrm{id}$. No appeal to finite-dimensionality is used.

A consequence we shall need twice: applying \eqref{eq:srdt-res} at the vector
$\xi=Tx$ gives $A_\varepsilon(Tx)=B\bigl(Tx+\varepsilon A_\varepsilon(Tx)\bigr)
=B\bigl(S(Tx)\bigr)=Bx$ by \eqref{eq:srdt-ST}, that is,
\begin{equation}\label{eq:srdt-AT}
  A_\varepsilon(Tx)=Bx\qquad(x\in V).
\end{equation}

\emph{(2b) The exact expansion of $F(\cdot,W)$ about $Z_W=TW$.}
Fix $W\in V$ and put $Z_W=TW$, so that $S Z_W=W$ by \eqref{eq:srdt-ST}. For any
$h\in V$, expanding the two $Z$-dependent terms of $F$ exactly,
\[
  \frac{|Z_W+h-W|^2}{2\varepsilon}
   =\frac{|Z_W-W|^2}{2\varepsilon}
     +\frac{\langle Z_W-W,h\rangle}{\varepsilon}
     +\frac{|h|^2}{2\varepsilon},
\]
\[
  \tfrac12\bigl\langle A_\varepsilon(Z_W+h),Z_W+h\bigr\rangle
   =\tfrac12\langle A_\varepsilon Z_W,Z_W\rangle
     +\langle A_\varepsilon Z_W,h\rangle
     +\tfrac12\langle A_\varepsilon h,h\rangle,
\]
where the second display uses the SELF-ADJOINTNESS OF $A_\varepsilon$ to
combine the two cross terms
$\tfrac12\langle A_\varepsilon Z_W,h\rangle
+\tfrac12\langle A_\varepsilon h,Z_W\rangle$ into
$\langle A_\varepsilon Z_W,h\rangle$. Subtracting both displays from
$\mathcal G(W)$ gives
\begin{equation}\label{eq:srdt-expand}
  F(Z_W+h,W)=F(Z_W,W)
    -\Bigl\langle\frac{Z_W-W}{\varepsilon}+A_\varepsilon Z_W,\,h\Bigr\rangle
    -\tfrac12\bigl\langle(\varepsilon^{-1}\mathrm{id}
      +A_\varepsilon)h,h\bigr\rangle .
\end{equation}
The linear term vanishes identically, because
\[
  \frac{Z_W-W}{\varepsilon}+A_\varepsilon Z_W
   =\frac{Z_W+\varepsilon A_\varepsilon Z_W-W}{\varepsilon}
   =\frac{S Z_W-W}{\varepsilon}=0 .
\]
Hence $F(Z_W+h,W)=F(Z_W,W)
-\tfrac12\langle(\varepsilon^{-1}\mathrm{id}+A_\varepsilon)h,h\rangle$ for
every $h$, and by \eqref{eq:srdt-pos} the subtracted quantity is $>0$ for
$h\ne0$. Therefore the supremum of $F(\cdot,W)$ is ATTAINED at $Z_W$ and
$Z_W$ is its UNIQUE maximizer:
\begin{equation}\label{eq:srdt-unique}
  F(Z,W)<F(Z_W,W)\quad(Z\ne Z_W),\qquad\sup_ZF(Z,W)=F(Z_W,W).
\end{equation}

\emph{(2c) The value at the maximizer.}
From $Z_W=TW=W-\varepsilon BW$ we get $Z_W-W=-\varepsilon BW$, so
$|Z_W-W|^2/(2\varepsilon)=\tfrac{\varepsilon}{2}|BW|^2$; and by
\eqref{eq:srdt-AT} with $x=W$ we have $A_\varepsilon Z_W=A_\varepsilon(TW)=BW$,
so
\[
  \langle A_\varepsilon Z_W,Z_W\rangle=\langle BW,TW\rangle
   =\langle BW,W-\varepsilon BW\rangle
   =\langle BW,W\rangle-\varepsilon|BW|^2 .
\]
Therefore
\begin{equation}\label{eq:srdt-duality}
  \sup_ZF(Z,W)=F(Z_W,W)
   =\mathcal G(W)-\frac{\varepsilon}{2}|BW|^2
     -\frac12\Bigl(\langle BW,W\rangle-\varepsilon|BW|^2\Bigr)
   =\mathcal G(W)-\tfrac12\langle BW,W\rangle=:\Phi(W).
\end{equation}

\emph{Step 3: $\Upsilon$ has a strict global maximum at $0$ with value $0$.}
By \eqref{eq:srdt-G3} we have $\Phi(W)\le-\sum_j|w_j|^4<0$ for $W\ne0$, while
$\Phi(0)=\sum_jG_j(0)-0=0$; so $\sup_W\Phi=0$, attained exactly at $W=0$. By
\eqref{eq:srdt-duality}, $\sup_{Z,W}F=\sup_W\Phi=0$. Hence $\Upsilon\le0$
everywhere, and $\Upsilon(0)=\sup_WF(0,W)\ge F(0,0)=\mathcal G(0)-0-0=0$, so
$\Upsilon(0)=0$.

Next, for each fixed $Z$ the supremum $\Upsilon(Z)=\sup_WF(Z,W)$ is attained:
$F(Z,\cdot)$ is continuous, and
$F(Z,W)\le NC-|Z-W|^2/(2\varepsilon)
-\tfrac12\langle A_\varepsilon Z,Z\rangle$, whose right-hand side tends to
$-\infty$ as $|W-Z|\to\infty$. Choose therefore $R>0$ so large that this bound
is $<F(Z,Z)$ whenever $|W-Z|>R$; then
$\sup_WF(Z,W)=\max_{\overline B(Z,R)}F(Z,\cdot)$, because every $W$ outside
$\overline B(Z,R)$ gives a value strictly below the value already attained at
$W=Z$, which lies inside the ball. The set $\overline B(Z,R)$ is closed,
bounded and nonempty, hence compact, so the maximum is attained (preamble).

Now suppose $\Upsilon(Z)=0$ for some $Z$. By the previous paragraph there is $W$
with $F(Z,W)=0$. Then $\Phi(W)=\sup_{Z'}F(Z',W)\ge F(Z,W)=0$, so $W=0$ by the
first paragraph of this step. Thus $F(Z,0)=0$. Now $Z_0=T\cdot0=0$, so the
expansion \eqref{eq:srdt-expand} at $W=0$ with $h=Z$ reads
\[
  F(Z,0)=F(0,0)
   -\tfrac12\bigl\langle(\varepsilon^{-1}\mathrm{id}+A_\varepsilon)Z,Z\bigr\rangle
   =-\tfrac12\bigl\langle(\varepsilon^{-1}\mathrm{id}
     +A_\varepsilon)Z,Z\bigr\rangle,
\]
using $F(0,0)=\mathcal G(0)-0-0=0$. Since $F(Z,0)=0$ this forces
$\langle(\varepsilon^{-1}\mathrm{id}+A_\varepsilon)Z,Z\rangle=0$, and
\eqref{eq:srdt-pos} then gives $Z=0$. This proves
\[
  \Upsilon(0)=0,\qquad\Upsilon(Z)<0\quad(Z\ne0).
\]

\emph{Step 4: $\Upsilon$ is $\lambda$-semiconvex and continuous.}
By part (3) of \noderef{SupConvolutionSemiconvex}, applied in $V$ to
$\mathcal G$ with growth constant $NC$, the function
$\mathcal G^\varepsilon$ is $(1/\varepsilon)$-semiconvex on $V$ in the sense of
\noderef{SemiconvexOn}, i.e.\ $\mathcal G^\varepsilon+|\cdot|^2/(2\varepsilon)$
is convex on $V$. Also
\[
  Z\mapsto-\tfrac12\langle A_\varepsilon Z,Z\rangle
    +\tfrac12\|A_\varepsilon\|\,|Z|^2
   =\tfrac12\bigl\langle(\|A_\varepsilon\|\mathrm{id}
     -A_\varepsilon)Z,Z\rangle
\]
is convex: for a self-adjoint $S$ with $\langle S\xi,\xi\rangle\ge0$ for all
$\xi$, and for $t_1,t_2\ge0$ with $t_1+t_2=1$, expanding and using
self-adjointness gives
\[
  \langle S(t_1\xi_1+t_2\xi_2),t_1\xi_1+t_2\xi_2\rangle
   -t_1\langle S\xi_1,\xi_1\rangle-t_2\langle S\xi_2,\xi_2\rangle
   =-t_1t_2\langle S(\xi_1-\xi_2),\xi_1-\xi_2\rangle\le0,
\]
and here $S=\|A_\varepsilon\|\mathrm{id}-A_\varepsilon$ is self-adjoint with
$\langle S\xi,\xi\rangle=\|A_\varepsilon\||\xi|^2
-\langle A_\varepsilon\xi,\xi\rangle\ge0$ by Cauchy--Schwarz and the definition
of the operator norm. Adding the two convex functions
$\mathcal G^\varepsilon+|\cdot|^2/(2\varepsilon)$ and
$-\tfrac12\langle A_\varepsilon\cdot,\cdot\rangle
+\tfrac12\|A_\varepsilon\||\cdot|^2$ shows that
$\Upsilon+\tfrac{\lambda}{2}|\cdot|^2$ is convex on $V$ with
$\lambda=\varepsilon^{-1}+\|A_\varepsilon\|>0$; that is, $\Upsilon$ is
$\lambda$-semiconvex on $V$.

Finally $\Upsilon+\tfrac\lambda2|\cdot|^2$ is a finite convex function on the
open set $V$, hence continuous (preamble); subtracting the continuous function
$\tfrac\lambda2|\cdot|^2$ shows that $\Upsilon$ is continuous on $V$.
\end{proof}
